\chapter{Path and Trajectory Planning}
\label{ch:path-trajectory}

Once the kinematic model of a wheeled mobile robot is known, the
planning problem becomes: find a path \( q(s) \) in the configuration
space that respects the kinematic constraints, and then a timing law
\( s(t) \) that makes it executable by the robot. This chapter
introduces the geometric model, the notion of differential flatness,
and a practical algorithm for path planning based on polynomial
interpolation of the flat outputs.

\section{Path vs Trajectory}
\label{sec:path-vs-trajectory}

We want to transfer a mobile robot from an initial configuration
\( q_i \) to a final configuration \( q_f \) without violating the
kinematic constraints on \( \dot{q} \). In general, we separate the
geometric path from the timing law:
\begin{align}
    q(s), \quad & s \in [s_i, s_f], \\
    s(t), \quad & t \in [t_i, t_f].
\end{align}
The path parameter \( s \) is usually chosen as the arc length.

\begin{itemize}
    \item \textbf{Path planning}: plan a geometric path \( q(s) \) in
          the configuration space.
    \item \textbf{Trajectory planning}: plan a time evolution
          \( q(t) \). This is usually done indirectly: first plan a
          path, then add a timing law.
\end{itemize}

A trajectory \( q(s(t)) \) must satisfy the kinematic constraints:
\begin{equation}
    \dot{q} = \frac{dq}{dt} = \frac{dq}{ds} \dot{s} = q' \dot{s},
\end{equation}
where \( q' \) denotes the derivative of \( q \) with respect to
\( s \). This is the \emph{tangent vector} of the geometric path in
the configuration space. If \( s \) is the arc length, then
\( \|q'\| = 1 \).

\paragraph{Example.}
Consider a robot with \( q = (x, y) \), and suppose we want to go from
\( (1, 0) \) to \( (0, 1) \). One possible choice for the geometric
path is
\[
q(s) = \begin{pmatrix} \cos s \\ \sin s \end{pmatrix},
\qquad s \in \left[0, \frac{\pi}{2}\right].
\]
This is exactly the arc length, because
\[
q'(s) = \begin{pmatrix} -\sin s \\ \cos s \end{pmatrix}
\implies \|q'(s)\| = \sqrt{\sin^2 s + \cos^2 s} = 1.
\]
Another parametrization is
\[
q(s) = \begin{pmatrix} \sqrt{1 - s^2} \\ s \end{pmatrix},
\qquad s \in [0, 1],
\]
but here \( s \) is not the arc length:
\[
q'(s) = \begin{pmatrix} -\frac{s}{\sqrt{1 - s^2}} \\ 1 \end{pmatrix}
\implies \|q'(s)\| = \frac{1}{\sqrt{1 - s^2}}.
\]

\section{The Geometric Model}
\label{sec:geometric-model}

The kinematic constraints are \( A(q) \dot{q} = 0 \). Substituting
\( \dot{q} = q' \dot{s} \), we get
\[
A(q) q' \dot{s} = 0.
\]
Assuming \( \dot{s} \neq 0 \), the tangent vector must satisfy
\[
q' \in \mathcal{N}(A(q)).
\]
Hence \( q' \) can be expressed as a linear combination of the basis
vectors \( g_1, \dots, g_m \), exactly as in the kinematic model. We
obtain the \textbf{geometric model}:
\begin{equation}
    q' = G(q) \tilde{u},
\end{equation}
which describes explicitly the admissible tangent vectors.
\( \tilde{u} \) is the vector of \textbf{geometric inputs}. Multiplying
both sides by \( \dot{s} \) gives
\[
\underbrace{u(t)}_{\text{pseudo-velocity}} = \tilde{u}(s) \dot{s}(t).
\]
After choosing a profile for the geometric inputs \( \tilde{u}(s) \),
we have determined a geometric path in the configuration space. Only
after defining also \( s(t) \) we obtain a particular trajectory.

\paragraph{Example: unicycle.}
For the unicycle model, the pure rolling constraint on the tangent
vector is
\[
\begin{pmatrix} \sin\theta & -\cos\theta & 0 \end{pmatrix} q' = 0.
\]
This means that the tangent vector \( q' \) must always be aligned with
the sagittal axis of the unicycle. Paths such as the one in
\cref{fig:inadmissible-path} are not admissible.

\begin{figure}[ht]
    \centering
    \includegraphics[width=0.4\linewidth]{images/inadmissible.eps}
    \caption{Geometric path with tangent vector not aligned with the
    sagittal axis.}
    \label{fig:inadmissible-path}
\end{figure}

\section{Differential Flatness}
\label{sec:differential-flatness}

\begin{definition}[Differential Flatness]
\label{def:flatness}
Consider a general nonlinear system
\[
\dot{x} = f(x) + G(x) u.
\]
The system is \textbf{differentially flat} if there exists a vector
\( y = \Psi(x) \), called the \textbf{flat outputs}, such that the
state and the control inputs can be expressed algebraically as a
function of \( y \) and its time derivatives up to a finite order:
\begin{align}
    x &= \Phi_x(y, \dot{y}, \ddot{y}, \dots, y^{(r_x)}), \\
    u &= \Phi_u(y, \dot{y}, \ddot{y}, \dots, y^{(r_u)}).
\end{align}
\end{definition}

The same definition holds if all derivatives are taken with respect to
the path parameter \( s \) instead of time, as in the geometric model.
Flatness is useful in control: if we define a trajectory for \( y \)
that is differentiable up to the order \( \max(r_x, r_u) \), we can
reconstruct the state evolution \( x \) and the required inputs \( u \)
through the reconstruction maps \( \Phi_x, \Phi_u \).

\paragraph{The unicycle is flat.}
For the unicycle model,
\begin{align}
    x' &= \tilde{v} \cos\theta, \\
    y' &= \tilde{v} \sin\theta, \\
    \theta' &= \tilde{\omega}.
\end{align}
The flat outputs are the Cartesian coordinates \( (x, y) \). If we
know the path of \( (x, y) \), we can infer the path of the orientation
\( \theta \) and the input profile that executes it. At any
\( s = s^* \), the orientation must be aligned with the direction of
motion, i.e.\ the vector \( (x'(s^*), y'(s^*)) \):
\begin{equation}
    \theta(s) = \text{atan2}(y'(s), x'(s)) + k\pi,
\end{equation}
where \( k \in \{0, 1\} \) defines the direction of motion
(forward or backward).

\begin{figure}[ht]
    \centering
    \includegraphics[width=0.6\linewidth]{images/flat_output_1.eps}
    \caption{Interpretation of the flat outputs \( (x,y) \) for the
    unicycle: the orientation is determined by the tangent to the
    Cartesian path.}
    \label{fig:flat-output-unicycle}
\end{figure}

The input \( \tilde{v} \) is the norm of the Cartesian velocity:
\begin{equation}
\label{eq:flat-v}
    \tilde{v} = \pm\sqrt{(x')^2 + (y')^2},
\end{equation}
and the angular velocity is
\begin{equation}
\label{eq:flat-omega}
    \tilde{\omega} = \frac{y'' x' - x'' y'}{(x')^2 + (y')^2}.
\end{equation}
Equations \eqref{eq:flat-v} and \eqref{eq:flat-omega} are the
reconstruction map for the input
\( \tilde{u} = \Phi_u(x', y', x'', y'') \). They have a singularity at
\( x' = y' = 0 \): atan2 is undefined and \( \tilde{\omega} \) has a
zero denominator. Intuitively, if the robot stops, we cannot know
whether it is steering in place or not.

\begin{nb}[Flatness and Chained-Form]
A driftless dynamic system is differentially flat if and only if it
can be transformed into chained-form. Since all the wheeled mobile
robots of interest admit a chained-form representation, they are also
flat.
\end{nb}

\paragraph{The bicycle is flat.}
For the bicycle model,
\begin{align}
    x' &= \tilde{v} \cos\theta, \\
    y' &= \tilde{v} \sin\theta, \\
    \theta' &= \frac{\tilde{v}}{l} \tan\phi, \\
    \phi' &= \tilde{\omega}.
\end{align}
The coordinates \( x, y, \theta \) are reconstructed as in the
unicycle case. The steering angle \( \phi \) is reconstructed from
\( \theta' \):
\begin{align}
    \theta(s) &= \text{atan2}(y', x') + k\pi, \\
    \theta'(s) &= \frac{x' y'' - y' x''}{(x')^2 + (y')^2},
\end{align}
hence
\begin{align}
    \phi(s) &= \arctan\!\left(\frac{l \theta'}{\tilde{v}}\right) \\
            &= \arctan\!\left(
               \pm \frac{l (x' y'' - y' x'')}
                         {\left((x')^2 + (y')^2\right)^{3/2}}
               \right).
\end{align}

\paragraph{Flat outputs for chained-form systems.}
Consider a 2-input, 3-state chained-form system:
\begin{align}
    z_1' &= \tilde{w}_1, \\
    z_2' &= \tilde{w}_2, \\
    z_3' &= z_2 \tilde{w}_1.
\end{align}
The flat outputs are always the first and the last variables:
\[
y = \begin{pmatrix} z_1 \\ z_3 \end{pmatrix}.
\]
The reconstruction maps are
\begin{align}
    \Phi_x(z_1', z_3') &= \frac{z_3'}{z_1'}, \\
    \Phi_u(z_1', z_1'', z_3', z_3'') &=
    \begin{cases}
        \tilde{w}_1 = z_1', \\[1ex]
        \tilde{w}_2 = \dfrac{z_1' z_3'' - z_1'' z_3'}{(z_1')^2}.
    \end{cases}
\end{align}

\section{Path Planning for the Unicycle}
\label{sec:path-planning-unicycle}

Any flat-output profile is automatically mapped to an admissible path.
The algorithm for path planning exploiting differential flatness is
the following:
\begin{enumerate}
    \item Compute the initial and final values \( y_i, y_f \) of the
          flat outputs from \( q_i, q_f \).
    \item Define \( y(s) \) by interpolating \( y_i, y_f \) with the
          appropriate boundary conditions for \( q \).
    \item Use the reconstruction formulas to find \( q(s) \) and
          \( \tilde{u}(s) \).
\end{enumerate}
The same algorithm applies to trajectory planning by interpolating
over \( t \) instead of \( s \).

Apply this to the unicycle. Use \( s \in [0, 1] \). From
\[
q_i = \begin{pmatrix} x_i \\ y_i \\ \theta_i \end{pmatrix},
\qquad
q_f = \begin{pmatrix} x_f \\ y_f \\ \theta_f \end{pmatrix},
\]
extract
\[
y_i = \begin{pmatrix} x_i \\ y_i \end{pmatrix},
\qquad
y_f = \begin{pmatrix} x_f \\ y_f \end{pmatrix},
\]
and interpolate with a cubic polynomial:
\begin{align}
    x(s) &= s^3 x_f - (s-1)^3 x_i + \alpha_x s^2 (s-1) + \beta_x s (s-1)^2, \\
    y(s) &= s^3 y_f - (s-1)^3 y_i + \alpha_y s^2 (s-1) + \beta_y s (s-1)^2.
\end{align}
This polynomial already satisfies the boundary conditions
\[
x(0) = x_i, \quad x(1) = x_f, \quad y(0) = y_i, \quad y(1) = y_f.
\]
The parameters \( \alpha_x, \beta_x, \alpha_y, \beta_y \) are chosen to
satisfy two additional boundary conditions from the kinematic
constraints: the tangent \( (x', y') \) must be aligned with
\( \theta_i \) at \( s = 0 \) and with \( \theta_f \) at \( s = 1 \).
Hence
\begin{align}
    \text{atan2}(y'(0), x'(0)) &= \theta_i, \\
    \text{atan2}(y'(1), x'(1)) &= \theta_f.
\end{align}
The initial and final Cartesian velocities must be
\begin{align}
    x'(0) &= k_i \cos\theta_i, & x'(1) &= k_f \cos\theta_f, \\
    y'(0) &= k_i \sin\theta_i, & y'(1) &= k_f \sin\theta_f,
\end{align}
with \( k_i, k_f \) arbitrary non-zero constants; for simplicity set
\( k_i = k_f = k \). Solving algebraically for
\( \alpha_x, \beta_x, \alpha_y, \beta_y \):
\begin{align}
    \alpha_x &= k \cos\theta_f - 3 x_f, & \beta_x &= k \cos\theta_i + 3 x_i, \\
    \alpha_y &= k \sin\theta_f - 3 y_f, & \beta_y &= k \sin\theta_i + 3 y_i.
\end{align}
Since \( \tilde{v}(0) = \tilde{v}(1) = k \), the parameter \( k \)
shapes the trajectory by affecting the length of the first and last
``linear'' segments of the path, as shown in
\cref{fig:k-shape}.

\begin{figure}[ht]
    \centering
    \includegraphics[width=0.8\linewidth]{images/k_shape.eps}
    \caption{Difference between two paths with the same interpolating
    points but different choice of \( k \).}
    \label{fig:k-shape}
\end{figure}

\section{Path Planning for a Chained-Form System}
\label{sec:path-planning-chained}

Apply the same algorithm to a 2-input, 3-state chained-form system.
The flat outputs are \( z_1, z_3 \), hence interpolate
\( (z_{1i}, z_{3i}) \) to \( (z_{1f}, z_{3f}) \). Interpolate \( z_1 \)
linearly and \( z_3 \) with a cubic polynomial:
\begin{align}
    z_1(s) &= z_{1i} + s (z_{1f} - z_{1i}), \\
    z_3(s) &= s^3 z_{3f} - (s-1)^3 z_{3i}
            + \alpha_3 s^2 (s-1) + \beta_3 s (s-1)^2,
\end{align}
which already satisfies the boundary conditions on \( z_1, z_3 \). The
additional boundary conditions for \( z_2 \) are
\[
z_{2i} = \frac{z_3'(0)}{z_1'(0)}, \qquad
z_{2f} = \frac{z_3'(1)}{z_1'(1)}.
\]
Solving for \( \alpha_3, \beta_3 \):
\begin{align}
    \alpha_3 &= z_{2f} (z_{1f} - z_{1i}) - 3 z_{3f}, \\
    \beta_3  &= z_{2i} (z_{1f} - z_{1i}) + 3 z_{3i}.
\end{align}
Once \( z_1(s), z_3(s) \) are known, reconstruct \( z_2(s),
\tilde{w}_1(s), \tilde{w}_2(s) \) through the reconstruction maps.
Then \( q(s), \tilde{v}(s), \tilde{\omega}(s) \) are obtained by
inverting the chained-form coordinate and input transformations.

A singularity arises when \( z_{1i} = z_{1f} \), because
\( z_1'(s) = z_{1f} - z_{1i} = 0 \), and \( z_{2i}, z_{2f} \) cannot
be inferred (division by zero). The fix is to set
\[
z_{1f} = z_{1i} + 2\pi,
\]
or to add an intermediate configuration
\( q_m = (x_m, y_m, \theta_m) \) with \( \theta_m \neq \theta_i \).

\section{Takeaways}
\label{sec:planning-takeaways}

\begin{enumerate}
    \item \textbf{Path planning} produces a geometric path
          \( q(s) \); \textbf{trajectory planning} adds a timing law
          \( s(t) \).
    \item The \textbf{geometric model} \( q' = G(q) \tilde{u} \) describes
          the admissible tangent vectors to the path.
    \item A driftless system is \textbf{differentially flat} if there
          exists a set of flat outputs from which state and inputs can
          be reconstructed algebraically. For the unicycle and the
          bicycle, the flat outputs are the Cartesian coordinates
          \( (x, y) \).
    \item Path planning via flatness reduces to polynomial interpolation
          of the flat outputs, with boundary conditions given by the
          initial and final configurations and the kinematic
          constraints.
    \item For chained-form systems, the flat outputs are the first and
          the last state variables; the reconstruction requires a
          careful handling of the singularity when the initial and
          final first coordinates coincide.
\end{enumerate}